\section{Conclusion}
\label{sec:conclusion}
This article treated a cross-silo federated job as a network tenant with a byte quota and asked when per-tier synchronisation pays. The closed form of its programme bounds the gain of any per-tier schedule over a uniform period by $\Psi$, measurable before any sparse round. On the main benchmark $\Psi$ is \num{auGateSpecialtyPsi}\% and \tierfed{} reduces to the uniform period; on a ResNet $\Psi$ is large and the tiers are split as predicted, \ifflag{auGateResnetWins}{and the split beats the best uniform period at matched bytes.}{but the split does not beat the best uniform period at matched bytes.} Phases still give tier structure something a uniform period cannot: a cap on the payload of any single round at an unchanged mean. Under client-level DP the schedule is part of the mechanism, and at 20 sites no releasing schedule comes near local training, so the quota and the privacy budget must be planned before training and checked against not federating at all.

\section*{Data, Code and Declarations}
Code, configuration, per-seed logs and the script that regenerates every table, figure and number are at \url{https://github.com/Abraheem13/tierfed}. All datasets are public and de-identified; no ethics approval was required. Generative artificial-intelligence tools assisted with code refactoring, figure layout and editing; the author verified all content, code, numbers and figures. An earlier version, ``Nested Federated Learning'', is superseded by this article (supplement, Section~S1). The author declares no competing interests and no funding.
